\section{Math}
\subsection{Math Formula Cheatsheet}
\textbf{Geometric Series}: \\
\textbullet{} $\sum_{i=a}^b c^i = \frac{c^{b+1} - c^a}{c - 1}$ $(c \neq 1)$ \\
\textbullet{} $\sum_{i=0}^n c^i = \frac{c^{n+1}-1}{c-1}$ $(c \neq 1)$ \\
\textbullet{} $\sum_{i=0}^\infty c^i = \frac{1}{1-c}$ $(|c| < 1)$

\textbf{Sums of Powers}: \\
\textbullet{} $\sum_{i=1}^n i = \frac{n(n+1)}{2}$ \\
\textbullet{} $\sum_{i=1}^n i^2 = \frac{n(n+1)(2n+1)}{6}$ \\
\textbullet{} $\sum_{i=1}^n i^3 = \frac{n^2(n+1)^2}{4}$ \\
\textbullet{} $\sum_{i=1}^n i^4 = \frac{n(n+1)(2n+1)(3n^2+3n-1)}{30}$

\textbf{Harmonic Sums}: \\
\textbullet{} $H_n = \sum_{i=1}^n \frac{1}{i}$ \\
\textbullet{} $\sum_{i=1}^n i H_i = \frac{n(n+1)}{2} H_n - \frac{n(n-1)}{4}$

\textbf{Combinatorial Sum}: \\
\textbullet{} $\sum_{k=0}^n \binom{r+k}{k} = \binom{r+n+1}{n}$

\subsection{Algebraic Identities}
\textbf{Squares}: $(a \pm b)^2 = a^2 \pm 2ab + b^2 \quad | \quad a^2 - b^2 = (a + b)(a - b)$ \\
\textbf{Cubes}: $(a \pm b)^3 = a^3 \pm 3a^2 b + 3ab^2 \pm b^3 \quad | \quad a^3 \pm b^3 = (a \pm b)(a^2 \mp ab + b^2)$ \\
\textbf{Trinomial}: $(a + b + c)^2 = a^2 + b^2 + c^2 + 2(ab + bc + ca)$ \\
\textbf{4th Power}: $a^4 + b^4 = (a^2 + b^2)^2 - 2a^2 b^2$ \\
\textbf{n-th Power}: $a^n - b^n = (a - b)(a^{n-1} + a^{n-2} b + \dots + b^{n-1})$

\subsection{Inequalities}
\textbf{AM-GM}: $\frac{x_1 + \dots + x_n}{n} \ge \sqrt[n]{x_1 \dots x_n}$ \\
\textbf{Cauchy-Schwarz}: $(\sum a_i^2)(\sum b_i^2) \ge (\sum a_i b_i)^2$ \\
\textbf{Triangle}: $||a| - |b|| \le |a \pm b| \le |a| + |b|$ \\
\textbf{Jensen}: $f(\frac{\sum x_i}{n}) \le \frac{\sum f(x_i)}{n}$ for convex $f$ \\
\textbf{Bernoulli}: $(1 + x)^n \ge 1 + nx$ for $n \ge 0, x > -1$ \\
\textbf{Chebyshev}: $\frac{1}{n} \sum a_i b_i \ge (\frac{1}{n} \sum a_i)(\frac{1}{n} \sum b_i)$ when sorted similarly \\
\textbf{Hölder}: $\sum |x_i y_i| \le (\sum |x_i|^p)^{1/p} (\sum |y_i|^q)^{1/q}$ where $1/p + 1/q = 1$

\subsection{Probability \& Expected Value}
$P(A \cup B) = P(A) + P(B) - P(A \cap B)$ \\
$P(A|B) = \frac{P(A \cap B)}{P(B)}$ (Bayes) \\
$E[X] = \sum x \cdot P(X = x)$ \\
$E[X + Y] = E[X] + E[Y]$ (Linearity) \\
$E[X \cdot Y] = E[X] \cdot E[Y]$ if independent \\
$Var(X) = E[X^2] - (E[X])^2$ \\
$Var(aX + b) = a^2 Var(X)$ \\
\textbf{Markov}: $P(X \ge a) \le \frac{E[X]}{a}$ \\
\textbf{Chebyshev}: $P(|X - \mu| \ge k\sigma) \le \frac{1}{k^2}$ \\
\textbf{Linearity of Expectation}: $E[\sum X_i] = \sum E[X_i]$ always \\
$Cov(X, Y) = E[XY] - E[X]E[Y]$
\subsection{Matrix Exponentiation}
\texttt{multiply(a,b)} multiplies matrices $O(N M P)$. \texttt{expo(t,p)} computes $t^p$ in $O(N^3 \log p)$. Skips zero entries in first matrix for speed. Define \texttt{mod} globally. Use for linear recurrences.
\begin{minted}{c++}
vector<vector<int>> multiply(
  vector<vector<int>>& a,
  vector<vector<int>>& b) {
  int n = sz(a), m = sz(a[0]), p = sz(b[0]);
  vector<vector<int>> result(n,
    vector<int>(p, 0));
  for (int i = 0; i < n; i++)
    for (int k = 0; k < m; k++) {
      if (a[i][k] == 0) continue;
      for (int j = 0; j < p; j++)
        result[i][j] = ((long long)result[i][j] +
          (long long)a[i][k] * b[k][j]) % mod;
    }
  return result;
}
vector<vector<int>> expo(
  vector<vector<int>> t, int p) {
  int n = t.size();
  vector<vector<int>> result(n, vector<int>(n, 0));
  for (int i = 0; i < n; i++) result[i][i] = 1;
  assert(p >= 0);
  while (p > 0) {
    if (p & 1) result = multiply(result, t);
    t = multiply(t, t);
    p >>= 1;
  }
  return result;
}
\end{minted}
\subsection{Basis Vector (XOR Basis)}
Maintains linear basis of vectors mod 2 (XOR space). \\
\texttt{insert(x)} -- returns true if x is independent. \\
\texttt{maxXor()} -- max XOR subset of inserted values. \\
\texttt{kth(k)} -- k-th smallest XOR (1-indexed). \\
\texttt{count\_less(x)} -- count of subsets with XOR $< x$. \\
\texttt{count\_leq(x)} -- count of subsets with XOR $\le x$. \\
\texttt{can(x)} -- checks if x is representable. \\
\texttt{reset()} -- clears the basis for reuse. \\
\texttt{jumpK(val, k)} -- k-th XOR value $\ge$ val (rank + k).
\begin{minted}{c++}
template<typename T = int, int LG = 31>
struct Basis {
  T basis[LG]; int sz;
  vector<int> who;
  Basis() { reset(); }
  void reset() {
    fill(basis, basis + LG, 0);
    sz = 0; who.clear();
  }
  bool insert(T mask, int idx = -1) {
    for (int i = LG - 1; i >= 0; --i) {
      if (mask >> i & 1) {
        if (basis[i]) mask ^= basis[i];
        else { basis[i] = mask; ++sz;
          who.push_back(idx); return true; }
      }
    }
    return false;
  }
  T maxXor() {
    T ans = 0;
    for (int i = LG - 1; i >= 0; --i)
      ans = max(ans, ans ^ basis[i]);
    return ans;
  }
  bool can(T x) {
    for (int i = LG - 1; i >= 0; --i) {
      if (x >> i & 1) {
        if (!basis[i]) return false;
        x ^= basis[i];
      }
    }
    return true;
  }
  T kth(T k) {
    if (k < 1 || (uint64_t)k > (1ULL << sz)) return -1;
    T mask = 0;
    int64_t tot = 1LL << sz;
    for (int i = LG - 1; i >= 0; --i) {
      if (basis[i]) {
        int64_t low = tot / 2;
        if ((low < k && (~mask >> i & 1)) ||
            (low >= k && (mask >> i & 1)))
          mask ^= basis[i];
        if (low < k) k -= low;
        tot /= 2;
      }
    }
    return mask;
  }
  T count_less(T x) {
    if (x < 0) return 0;
    T ret = 0, mask = 0;
    uint64_t cnt = 1ULL << sz;
    for (int i = LG - 1; i >= 0; i--) {
      if (basis[i]) {
        if (x >> i & 1) {
          ret += (cnt >> 1);
          if (~mask >> i & 1) mask ^= basis[i];
        } else {
          if (mask >> i & 1) mask ^= basis[i];
        }
        cnt >>= 1;
      } else {
        if ((x ^ mask) >> i & 1) {
          if (x >> i & 1) return ret + cnt;
          else return ret;
        }
      }
    }
    return ret;
  }
  T count_leq(T x) { return count_less(x) + (can(x) ? 1 : 0); }
  T jumpK(T val, T k) {
    return kth(count_leq(val) + k);
  }
};
\end{minted}
\subsection{Linear Combination Checking}
Given array \texttt{arr} and integer N, check if N is
a linear combination of arr's elements:
$N = a \cdot x + b \cdot y + c \cdot z + \ldots$

Ans: $N \bmod \text{GCD}(\text{all elements}) == 0$.
\textbf{Note:} A sum $S$ can be formed if $S \pmod{\gcd} == 0$.

\subsection{FFT (Fast Fourier Transform)}
Convolution for integer polynomials (double-based). \texttt{multiply(a,b)} returns $a * b$ as \texttt{vector<long long>}.
Prefer NTT when modulus supports it (no precision issues). Complexity $O(N \log N)$ where $N$ is next power of 2.
\textbf{Note:} Array sizes must be a power of 2. Pad arrays with 0s before calling.
\begin{minted}{c++}
const double PI = acos(-1);
struct base {
  double a, b;
  base(double a = 0, double b = 0) : a(a), b(b) {}
  const base operator+(const base &c) const
    { return base(a + c.a, b + c.b); }
  const base operator-(const base &c) const
    { return base(a - c.a, b - c.b); }
  const base operator*(const base &c) const
    { return base(a*c.a - b*c.b, a*c.b + b*c.a); }
};
void fft(vector<base> &p, bool inv = 0) {
  int n = p.size(), i = 0;
  for (int j = 1; j < n - 1; ++j) {
    for (int k = n >> 1; k > (i ^= k); k >>= 1);
    if (j < i) swap(p[i], p[j]);
  }
  for (int l = 1; (l << 1) <= n; l <<= 1) {
    double ang = 2 * PI / (l << 1);
    base wn(cos(ang), (inv ? 1 : -1) * sin(ang)), w;
    for (int i = 0; i < n; i += l << 1) {
      w = base(1, 0);
      for (int j = i, k = i + l; j < k;
           ++j, w = w * wn) {
        base t = w * p[j + l];
        p[j + l] = p[j] - t;
        p[j] = p[j] + t;
      }
    }
  }
  if (inv)
    for (int i = 0; i < n; ++i)
      p[i].a /= n, p[i].b /= n;
}
vector<long long> multiply(vector<long long> &a,
                            vector<long long> &b) {
  if (a.empty() || b.empty()) return {};
  int n = a.size(), m = b.size(), t = n + m - 1, sz = 1;
  while (sz < t) sz <<= 1;
  vector<base> x(sz), y(sz);
  for (int i = 0; i < sz; ++i) {
    x[i] = i < n ? base(a[i], 0) : base(0, 0);
    y[i] = i < m ? base(b[i], 0) : base(0, 0);
  }
  fft(x), fft(y);
  for (int i = 0; i < sz; ++i) x[i] = x[i] * y[i];
  fft(x, 1);
  vector<long long> ret(sz);
  for (int i = 0; i < sz; ++i)
    ret[i] = (long long)round(x[i].a);
  while (ret.size() > 1 && ret.back() == 0)
    ret.pop_back();
  return ret;
}
\end{minted}
\subsection{NTT (Number Theoretic Transform)}
Modular convolution ($\bmod 998244353$, primitive root $= 62$).
Inputs must be in $[0, \text{mod})$ before calling. \texttt{conv(a,b)} returns polynomial product modulo MOD.
Replace MOD and ROOT for different moduli ($\text{mod} = c \cdot 2^k + 1$).
\textbf{Note:} The modulo must be of the form $c \cdot 2^k + 1$ (like $998244353$).
\begin{minted}{c++}
const int MOD = 998244353, ROOT = 62;
int modpow(int a, int e) {
  int r = 1;
  while (e) {
    if (e & 1) r = (long long)r * a % MOD;
    a = (long long)a * a % MOD;
    e >>= 1;
  }
  return r;
}
void ntt(vector<int> &a) {
  int n = a.size(), L = 31 - __builtin_clz(n);
  static vector<int> rt(2, 1);
  for (static int k = 2, s = 2; k < n; k *= 2, s++) {
    rt.resize(n);
    int z[] = {1, modpow(ROOT, MOD >> s)};
    for (int i = k; i < 2 * k; i++)
      rt[i] = (long long)rt[i / 2] * z[i & 1] % MOD;
  }
  vector<int> rev(n);
  for (int i = 0; i < n; i++)
    rev[i] = (rev[i / 2] | (i & 1) << L) / 2;
  for (int i = 0; i < n; i++)
    if (i < rev[i]) swap(a[i], a[rev[i]]);
  for (int k = 1; k < n; k *= 2)
    for (int i = 0; i < n; i += 2 * k)
      for (int j = 0; j < k; j++) {
        int z = (long long)rt[j + k] * a[i + j + k] % MOD;
        a[i + j + k] = a[i + j] - z +
          (z > a[i + j] ? MOD : 0);
        a[i + j] += a[i + j] + z >= MOD ? z - MOD : z;
      }
}
vector<int> conv(const vector<int> &a,
                 const vector<int> &b) {
  if (a.empty() || b.empty()) return {};
  int s = a.size() + b.size() - 1,
      n = 1 << (32 - __builtin_clz(s));
  int inv = modpow(n, MOD - 2);
  vector<int> L(a), R(b), out(n);
  L.resize(n); R.resize(n);
  ntt(L); ntt(R);
  for (int i = 0; i < n; i++)
    out[(n - i) & (n - 1)] = (long long)L[i] * R[i] %
      MOD * inv % MOD;
  ntt(out);
  return {out.begin(), out.begin() + s};
}
\end{minted}
\subsection{FWHT}
Convolution where $c[k] = \sum_{i \oplus j = k} a[i] \cdot b[j]$.
\texttt{type}: 0=OR, 1=AND, 2=XOR. \texttt{inv=true} for inverse. \texttt{modpow(n, MOD-2)} for inverse normalization (type 2 only).
\textbf{Note:} Array size must be a power of 2.
\begin{minted}{c++}
void fwht(vector<int> &a, bool inv, int type) {
  int n = a.size();
  for (int len = 1; len < n; len <<= 1)
    for (int i = 0; i < n; i += 2 * len)
      for (int j = 0; j < len; j++) {
        int u = a[i + j], v = a[i + len + j];
        if (type == 2) {
          a[i + j] = (u + v) % MOD;
          a[i + len + j] = (u - v + MOD) % MOD;
        } else if (type == 1) {
          a[i+j] = inv ? (u-v+MOD)%MOD : (u+v)%MOD;
        } else {
          a[i+len+j] = inv ? (v-u+MOD)%MOD : (u+v)%MOD;
        }
      }
  if (inv && type == 2) {
    int inv_n = modpow(n, MOD - 2);
    for (auto &x : a) x = (ll)x * inv_n % MOD;
  }
}
// conv(a, b) = for type=0: c[k]=sum_{i|j=k} a[i]*b[j]
// for type=1: c[k]=sum_{i&j=k} a[i]*b[j]
// for type=2: c[k]=sum_{i^j=k} a[i]*b[j]
\end{minted}
\subsection{Subset Convolution}
Computes $h[mask] = \sum_{sub \subseteq mask} f[sub] \cdot g[mask \setminus sub]$ where $sub$ and $mask \setminus sub$ are \textbf{disjoint}. Unlike FWHT OR-convolution (which allows overlap), subset convolution enforces disjointness via an extra dimension for population count. $O(B^2 \cdot 2^B)$.
\textbf{When to use:} Set partition DP, graph matching on small graphs, counting $k$-colorings, any DP where two disjoint subsets are combined.
\begin{minted}{c++}
vector<int> subset_conv(const vector<int>& f,
                        const vector<int>& g, int B) {
  int M = 1 << B;
  vector<vector<int>> F(B+1, vector<int>(M));
  vector<vector<int>> G(B+1, vector<int>(M));
  vector<vector<int>> H(B+1, vector<int>(M));
  for (int m = 0; m < M; m++) {
    int p = __builtin_popcount(m);
    F[p][m] = f[m]; G[p][m] = g[m];
  }
  // Zeta transform (SOS) on each rank
  for (int k = 0; k <= B; k++)
    for (int b = 0; b < B; b++)
      for (int m = 0; m < M; m++)
        if (m >> b & 1) {
          F[k][m] = (F[k][m] + F[k][m ^ (1<<b)]) % MOD;
          G[k][m] = (G[k][m] + G[k][m ^ (1<<b)]) % MOD;
        }
  // Pointwise multiply over ranks
  for (int m = 0; m < M; m++)
    for (int k = 0; k <= B; k++)
      for (int i = 0; i <= k; i++)
        H[k][m] = (H[k][m] +
          (ll)F[i][m] * G[k-i][m]) % MOD;
  // Mobius transform (inverse zeta)
  for (int k = 0; k <= B; k++)
    for (int b = 0; b < B; b++)
      for (int m = 0; m < M; m++)
        if (m >> b & 1)
          H[k][m] = (H[k][m] - H[k][m ^ (1<<b)]
                     + MOD) % MOD;
  vector<int> h(M);
  for (int m = 0; m < M; m++)
    h[m] = H[__builtin_popcount(m)][m];
  return h;
}
\end{minted}

\subsection{Berlekamp-Massey}
\textbf{What it is:} A "black box" algorithm that looks at a sequence of numbers (like the first few terms of a DP) and automatically discovers the hidden formula (linear recurrence) generating them. E.g., for Fibonacci, it discovers $a_n = a_{n-1} + a_{n-2}$.

\textbf{How to use:} Compute the first $\approx 100$ terms of your sequence using naive DP and pass them as a vector. It returns the recurrence coefficients $C$. You can then use Matrix Exponentiation (or Kitamasa) with $C$ to find the $K$-th term for a massive $K$.

\textbf{Note:} Needs at least $2L$ terms to correctly find a recurrence of length $L$.
\begin{minted}{c++}
vector<int> berlekampMassey(vector<int> x) {
  int n = x.size(), L = 0, m = 1;
  ll b = 1;
  vector<int> C(n), B(n);
  C[0] = B[0] = 1;
  for (int i = 0; i < n; i++, m++) {
    ll d = x[i] % MOD;
    for (int j = 1; j <= L; j++)
      d = (d + (ll)C[j] * x[i - j]) % MOD;
    if (!d) continue;
    vector<int> T = C;
    ll c = d * modpow(b, MOD - 2) % MOD;
    for (int j = m; j < n; j++)
      C[j] = (C[j] - c * B[j - m] + MOD) % MOD;
    if (2 * L <= i) {
      L = i - L + 1; B = T;
      b = d; m = 0;
    }
  }
  C.resize(L + 1); C.erase(C.begin());
  for (auto &i : C) i = (MOD - i) % MOD;
  return C;
}
// Kitamasa: compute k-th term of linear recurrence
// using C (from BM) and first terms.
// Complexity O(L^2 log k).
\end{minted}
